\section{Introducción}
\label{sec:introduccion}

Los dos reportes previos de este proyecto establecieron que un sistema de agentes
especialistas construido exclusivamente con memorias asociativas entrópicas desarrolla
el directorio transactivo de Wegner \cite{wegner1987transactive} por interacción
acumulativa: primero con tres agentes (precisión madura 98.75\% en tres categorías de
ETH-80) y después con ocho, cubriendo el conjunto ETH-80 \cite{leibe2003eth80} completo
(precisión temprana 88.0\%, madura 93.2\%, hemisferio visual con 75.0\% de
enrutamiento y cero falsos aceptados).
En ambos casos la evaluación fue \emph{por lotes}: bancos de consultas escritas y
archivos de imagen procesados fuera de línea.

Este reporte aborda la extensión natural que esos resultados dejan pendiente, y que
en la literatura se conoce como la construcción de un \textbf{sistema multimodal}
\cite{baltrusaitis2019multimodal}: un sistema que percibe y coordina información
proveniente de más de una modalidad sensorial.
La pregunta que organiza el trabajo es:

\begin{quote}
¿Puede la maquinaria transactiva sobre memorias asociativas entrópicas operar en
\emph{tiempo real} con entrada sensorial continua ---voz hablada y video de una cámara
web--- manteniendo el principio de que \emph{todas} las decisiones son operaciones de
memoria? Y cuando la memoria hetero-asociativa recupera sin pista en el dominio
destino, ¿cómo se pasa de la masa indeterminada que devuelve la proyección
---el \emph{problema de la pista faltante} \cite{morales2025missing}--- a un objeto
definido?
\end{quote}

La respuesta empírica es afirmativa, con una condición de diseño que este reporte
documenta: las \emph{traducciones} entre modalidades viven en los bordes del sistema
(voz$\to$texto mediante Whisper \cite{radford2023whisper}; pixeles$\to$latente mediante
el encoder ResNet18 \cite{he2016resnet} ya existente), mientras que el núcleo
---reconocer, asignar, redirigir, rechazar y recuperar--- permanece intacto: son las
mismas operaciones de las memorias W-EAM y EHAM
\cite{pineda2021entropic,pineda2022weighted,morales2024hetero} evaluadas en los
reportes anteriores, consultadas ahora en modo de solo lectura por una ventana nativa
de tres hilos.

\subsection*{Contribuciones principales}

Este trabajo aporta seis contribuciones concretas respecto a la versión de ocho
agentes:

\begin{enumerate}
    \item \textbf{Canal de voz continuo.}
    Escucha permanente con detección de actividad de voz (Silero VAD
    \cite{silero2021vad}), subtítulos parciales en vivo y enrutamiento automático al
    detectar la pausa del hablante ($\sim$600 ms de silencio).
    La transcripción usa Whisper \emph{small} \cite{radford2023whisper} en GPU
    (vía faster-whisper \cite{systran2023fasterwhisper}) en modo \texttt{translate},
    de modo que el usuario puede hablar en español o inglés y la memoria recibe pistas
    en inglés, el idioma de su vocabulario fastText.
    Latencia medida extremo a extremo (fin de habla $\to$ decisión de enrutamiento):
    $\approx$0.6 s.

    \item \textbf{Canal de visión con etiqueta viva.}
    El cuadro de la cámara se analiza $\sim$4 veces por segundo (54 ms por análisis)
    con suavizado por mayoría, mostrando sobre el video la clase reconocida por el
    directorio visual del agente de entrada o un \emph{rechazo honesto} con el
    candidato más cercano.
    La interfaz muestra además ``lo que ve el ojo'': el recorte de $128\times128$
    exacto que recibe el encoder, que convierte los fallos de encuadre, enfoque o
    reflejo en evidencia visible al instante.

    \item \textbf{Los dos modos transactivos, en vivo y visibles.}
    La voz opera en modo \emph{temprano} (difusión del TME: los ocho agentes puntúan y
    el grupo asigna, mientras un directorio de sesión aprende de cada frase) o
    \emph{maduro} (protocolo punto a punto: el agente de entrada consulta su directorio
    entrenado con lectura B1 y redirige al especialista o rechaza, con el TME apagado).
    El panel muestra la redirección como relación entre agentes
    (\textsc{entrada} $\to$ \textsc{destino}), haciendo observable el mecanismo de
    Wegner en tiempo real.

    \item \textbf{Política de color ``reintentar con lentes''.}
    La cadena física impresora$+$cámara desplaza los momentos de color del estímulo y
    colapsa el enrutamiento visual.
    Se caracterizó el fenómeno, se demostró que un detector por distancia de momentos
    no separa las entradas limpias de las degradadas (las distribuciones se traslapan),
    y se adoptó una política sin umbral: intento crudo primero y renormalización de
    momentos RGB hacia los anclajes de ETH-80 \emph{solo} si el directorio rechaza.
    El rescate triplica el reconocimiento bajo degradación fuerte (de $\sim$0 a 15/24)
    sin dañar las entradas limpias.

    \item \textbf{Recuperación con pista faltante: los tres métodos del artículo.}
    Se implementaron \emph{random samples} (RS), \emph{sample-and-test} (ST) y
    \emph{sample-and-search} (SS) de Morales \& Pineda \cite{morales2025missing} como
    métodos seleccionables de la memoria hetero-asociativa, compartiendo la misma
    proyección AND-ponderada y la misma prueba de retro-proyección.
    Se verificó su jerarquía empírica en este dominio y se documentó que el
    mecanismo de recuperación oficial del sistema es exactamente SS, el mejor de los
    tres.

    \item \textbf{Prototipo emergente vs.\ objeto definido.}
    Ambas interfaces (la aplicación de trazado del experimento y la ventana en tiempo
    real) exhiben ahora el contraste que motiva el problema de la pista faltante: el
    \emph{plano proyectado} como masa indeterminada, su \emph{prototipo emergente}
    (la lectura determinista columna a columna, que rompe la consistencia conjunta) y
    el \emph{objeto definido} que produce la búsqueda con prueba, comparados con una
    métrica común: la distancia de retro-proyección a la pista original.
\end{enumerate}

\subsection*{Alcance}

El experimento base permanece intacto: la extensión vive en un módulo aislado
(\texttt{realtime\_lab/}) que carga los modelos entrenados en modo de solo lectura y
mantiene su propio directorio de sesión en memoria.
Ningún número de los reportes anteriores cambia; los resultados nuevos son las
latencias, los diagnósticos de la cadena física de percepción y la caracterización de
los métodos de recuperación con pista faltante.

\subsection*{Estructura del documento}

La Sección~\ref{sec:marco_teorico} resume los fundamentos y formaliza el problema de
la pista faltante y sus tres métodos de solución.
La Sección~\ref{sec:arquitectura} describe la arquitectura multimodal en tiempo real.
La Sección~\ref{sec:metodologia} detalla la metodología de verificación.
La Sección~\ref{sec:resultados} reporta latencias, diagnósticos físicos y la
caracterización empírica de los métodos de recuperación.
La Sección~\ref{sec:discusion} discute los hallazgos y la
Sección~\ref{sec:conclusiones} concluye.
